% Figure: an exponential decay plotted on linear axes (left) and on
% semilog axes (right). The same data, the same physics; the semilog
% version reads the time constant as a slope and the initial value
% as the intercept.
\begin{figure}[htbp]
\centering
\resizebox{\linewidth}{!}{%
\begin{tikzpicture}[scale=1.0,
  ax/.style={->, draw=engblack, thick},
  tick/.style={draw=engblack!60, thin},
  curve/.style={draw=engred, thick},
  fitline/.style={draw=engblue, thick, dashed},
  pt/.style={circle, fill=engred, inner sep=1.2pt},
  ann/.style={font=\footnotesize, align=center},
]

% --- Left: linear ---
\begin{scope}[xshift=0cm]
  \draw[ax] (0, 0) -- (5.8, 0) node[below right, ann] {$t$ (s)};
  \draw[ax] (0, 0) -- (0, 4.2) node[left, ann] {$V$ (V)};
  \foreach \x in {1, 2, 3, 4, 5} {
    \draw[tick] (\x, -0.05) -- (\x, 0.05);
    \node[ann, below, yshift=-2pt] at (\x, 0) {\x};
  }
  \foreach \y in {1, 2, 3, 4} {
    \draw[tick] (-0.05, \y) -- (0.05, \y);
    \node[ann, left, xshift=-2pt] at (0, \y) {\y};
  }
  % V(t) = 4 e^{-t/1.5}, scaled to fit
  \draw[curve, domain=0:5.5, samples=40, smooth]
    plot (\x, {4*exp(-\x/1.5)});
  % Sampled points
  \foreach \x/\y in {0/4.0, 0.5/2.87, 1.0/2.05, 1.5/1.47,
    2.0/1.05, 2.5/0.75, 3.0/0.54, 3.5/0.39, 4.0/0.28, 4.5/0.20, 5.0/0.14} {
    \node[pt] at (\x, \y) {};
  }
  \node[ann, anchor=south] at (3.0, 3.5) {(a) linear axes};
\end{scope}

% --- Right: semilog (vertical log) ---
\begin{scope}[xshift=8cm]
  \draw[ax] (0, 0) -- (5.8, 0) node[below right, ann] {$t$ (s)};
  \draw[ax] (0, -0.5) -- (0, 4.2) node[left, ann] {$\log_{10} V$};
  \foreach \x in {1, 2, 3, 4, 5} {
    \draw[tick] (\x, -0.05) -- (\x, 0.05);
    \node[ann, below, yshift=-2pt] at (\x, 0) {\x};
  }
  % Major decades; log V scaled so 1 unit = one decade
  \foreach \y/\yl in {0/0, 1/1, 2/2, 3/3, 4/4} {
    \draw[tick] (-0.05, \y) -- (0.05, \y);
    \node[ann, left, xshift=-2pt] at (0, \y) {$10^{\yl}$};
  }
  % Plot log10(4 e^{-t/1.5}) shifted vertically by +2 so it stays
  % visible: y_axis = 2 + log10(4) + (-t/1.5)*log10(e)
  % = 2.602 - 0.2895 * t
  \draw[fitline, domain=0:5.5, samples=2]
    plot (\x, {2.602 - 0.2895*\x});
  \foreach \x in {0, 0.5, 1.0, 1.5, 2.0, 2.5, 3.0, 3.5, 4.0, 4.5, 5.0} {
    \pgfmathsetmacro{\yy}{2.602 - 0.2895*\x}
    \node[pt] at (\x, \yy) {};
  }
  \node[ann, anchor=south] at (3.0, 3.5) {(b) semilog axes};
  \node[ann, anchor=west, draw=engblue, fill=white, inner sep=2pt]
    at (1.4, 3.0)
    {slope $= -\log_{10}(e)/\tau$\\$\tau$ from slope};
\end{scope}

\end{tikzpicture}%
}%
\caption[An exponential decay on linear and semilog axes]{The same
exponential decay $V(t) = V_{0} e^{-t/\tau}$ with $V_{0} =
4\,\si{\volt}$ and $\tau = 1.5\,\si{\second}$, plotted on linear
axes (left) and on a vertical-log axis (right). On linear axes the
curve is the recognisable convex decay that the eye reads as ``fast
at first, then slow''. On semilog axes the same data is a straight
line, with intercept $\log_{10} V_{0}$ and slope $-\log_{10}(e)/\tau
\approx -0.290\,\si{\per\second}$. The engineer who fits a straight
line to the semilog data recovers $\tau$ from the slope in one step
($\tau = -\log_{10}(e) / \text{slope}$) and $V_{0}$ from the
intercept in another. A straight-line fit on the linear panel
recovers nothing.}
\label{fig:vol02:ch02:semilog-decay}
\end{figure}
